\documentclass[runningheads]{llncs}

\usepackage[T1]{fontenc}
\usepackage[utf8]{inputenc}
\usepackage{graphicx}
\usepackage{amsmath,amssymb}
\usepackage{booktabs}
\usepackage{multirow}
\usepackage{xcolor}
\usepackage{hyperref}
\usepackage{enumitem}
\usepackage{algorithm}
\usepackage{algpseudocode}
\setlist{nosep}

\hypersetup{colorlinks=true, linkcolor=blue!60!black,
            citecolor=blue!60!black, urlcolor=blue!60!black}

\begin{document}

\title{Scalable Quantum Neighborhood Search for the Permutation
       Flow Shop Problem via Windowed QUBO Decomposition}

\titlerunning{Scalable Quantum Neighborhoods for PFSP}

\author{Remigiusz Wojew\'{o}dzki\inst{1} \and
        Wojciech Bo\.{z}ejko\inst{1}}
\authorrunning{R. Wojew\'{o}dzki, W. Bo\.{z}ejko}

\institute{
  Department of Control Systems and Mechatronics,\\
  Wroc\l{}aw University of Science and Technology, Poland\\
  \email{\{remigiusz.wojewodzki,wojciech.bozejko\}@pwr.edu.pl}
}

\maketitle

\begin{abstract}
We extend the quantum neighborhood framework for the Permutation Flow
Shop Scheduling Problem (PFSP) by adding a third neighborhood class,
\emph{quantum QUBO enhanced}, which runs on significantly larger
benchmark instances than the original D-Wave formulations allowed.
The prior framework had eight neighborhoods: four classical
(Adjacent, Fibonacci, Dynasearch, Motzkin) and four quantum QUBO
variants limited to $n = 20$ jobs on D-Wave hardware.
The enhanced class fixes two problems: an unnecessary delta-filtering
heuristic in linearly-scaling QUBO structures (Adjacent and Fibonacci)
is removed, extending QPU applicability to $n \leq 200$ and $n \leq 500$;
and a windowed QUBO decomposition is added for quadratically-scaling
structures (Dynasearch and Motzkin), pushing feasibility to $n \leq 50$.
Both metaheuristic frameworks --- Iterated Local Search with tabu
diversification via a MushroomList elite pool, and Simulated Annealing
with adaptive reheating --- use a block-property accelerator that prunes
dominated swap candidates before each QUBO submission, cutting the
effective variable count by 30--60\% on Taillard instances.
All twelve neighborhoods are evaluated on standard Taillard benchmarks
across six time budgets (100--10\,000\,ms).
\keywords{Flow shop scheduling \and QUBO \and Quantum annealing \and
          D-Wave \and Neighborhood search \and Windowed decomposition \and
          MushroomList}
\end{abstract}

% ---------------------------------------------------------------
\section{Introduction}

The Permutation Flow Shop Scheduling Problem (PFSP) asks for a
permutation of $n$ jobs on $m$ machines in series that minimizes
makespan $C_{\max}$. The problem is NP-hard for $m \geq 3$~\cite{ref_garey},
and for large benchmark instances the best known solutions are
produced by neighborhood-based metaheuristics rather than exact methods.
The choice of neighborhood structure — which candidate permutations to
examine at each step — is therefore a primary driver of solution quality.

In~\cite{WojewodzkiBozejko2026} we formulated four classical PFSP
neighborhoods as Quadratic Unconstrained Binary Optimization (QUBO)
problems and executed them on D-Wave quantum annealing hardware, providing
the first systematic comparison of classical and quantum neighborhood
search under identical wall-clock time budgets and metaheuristic frameworks.
The experimental scope was restricted to Taillard instances with $n = 20$
jobs, imposed jointly by the variable-count constraints of the QUBO
formulations and the effective embedding capacity of the D-Wave Advantage
QPU (approximately 180 effective variables for dense QUBO graphs).

This paper addresses both restrictions directly. Two properties in the
original formulations block scalability:

\begin{enumerate}
  \item \textbf{Unnecessary delta filtering.}
        The original implementation discarded swap candidates with
        non-negative makespan delta ($\delta_k \geq 0$) before QUBO
        assembly --- a workaround for the classical simulator that
        is counterproductive on real QPU hardware.
        The QPU assigns $x_k = 0$ to non-improving candidates on its own,
        but the filter discards information about joint penalty structure
        that helps the QPU avoid degenerate solutions.
        For Adjacent and Fibonacci, where $K = n - 1$ regardless of
        filtering, removing it costs nothing in embedding overhead.

  \item \textbf{Variable-count ceiling for $O(n^2)$ neighborhoods.}
        Dynasearch and Motzkin have $K = O(n^2)$ QUBO variables,
        exceeding QPU capacity for $n > 20$.
        A windowed decomposition partitions the permutation into
        overlapping windows of size $w \leq 19$, solves a QUBO per
        window, and merges results via greedy conflict resolution.
\end{enumerate}

Together these give twelve neighborhoods in three classes --- classical,
quantum QUBO, and quantum QUBO enhanced --- covering the full Taillard
benchmark range for the first time.

\paragraph{Contributions.}
\begin{enumerate}
  \item A windowed QUBO decomposition for Dynasearch and Motzkin
        neighborhoods that makes QPU execution feasible for $n \leq 50$,
        with automatic window sizing from hardware capacity.
  \item Delta-filter-free QUBO formulations for Adjacent ($n \leq 200$)
        and Fibonacci ($n \leq 500$) on D-Wave Advantage.
  \item A block-property accelerator (Smutnicki~\cite{ref_smutnicki})
        applied to all twelve neighborhoods, cutting the QUBO variable
        count by 30--60\% before QPU submission.
  \item A MushroomList elite-pool for ILS that replaces random restarts
        with double-bridge perturbations of the $k = 10$ best solutions
        seen during search.
  \item The first empirical comparison of all three neighborhood classes
        across $n \in \{20, 50, 100, 200, 500\}$ under identical
        metaheuristic frameworks and time budgets.
\end{enumerate}

% ---------------------------------------------------------------
\section{Problem Definition}
\label{sec:problem}

\paragraph{Permutation Flow Shop Problem.}
Let $P = [p_{k,j}]_{m \times n}$ denote processing times, where machine
index $k \in \{1,\ldots,m\}$ and job index $j \in \{1,\ldots,n\}$.
A permutation $\pi \in S_n$ specifies the processing order on all machines.
Completion time $C_{k,j}(\pi)$ satisfies
\begin{align}
  C_{1,1} &= p_{1,\pi(1)}, \quad
  C_{k,1} = C_{k-1,1} + p_{k,\pi(1)}, \quad
  C_{1,j} = C_{1,j-1} + p_{1,\pi(j)}, \\
  C_{k,j} &= \max\!\bigl\{C_{k,j-1},\, C_{k-1,j}\bigr\} + p_{k,\pi(j)}.
\end{align}
The objective is $\pi^* = \arg\min_\pi C_{m,n}(\pi) =: C_{\max}(\pi)$.

\paragraph{Head--Tail decomposition.}
Following~\cite{wodecki04}, the makespan delta for inserting job
$\pi(i)$ at position $j$ is computed in $O(1)$ using precomputed
head matrices $H[k][j]$ (earliest completion time of job at position
$j$ on machine $k$ from left) and tail matrices $T[k][j]$ (latest
completion time contribution from right). Both matrices are built in
$O(mn)$ before each local search iteration and reused for all
candidate evaluations within that iteration.

\paragraph{Metaheuristics.}
We evaluate all neighborhoods under two frameworks.

\emph{Iterated Local Search (ILS)~\cite{lourenco2003iterated}}
with a tabu list of tenure $\tau = 10$ and aspiration criterion
(accept a tabu move if it strictly improves the global best $C_{\max}^*$).
When no non-tabu move is available and aspiration is not met, the
search diversifies via the MushroomList mechanism described in
Section~\ref{sec:mushroom}.

\emph{Simulated Annealing (SA)~\cite{kirkpatrick1983}}
with initial temperature $T_0 = 1000$, cooling rate $\alpha = 0.95$,
and adaptive reheating: when no improvement occurs for 500\,ms the
temperature is reset to $\min(T_0,\, T \cdot r)$ with reheat factor
$r = 1.5$.

\subsection{MushroomList Elite Diversification}
\label{sec:mushroom}

Uniform random restarts reinitialize from an arbitrary point and throw
away everything learned up to that moment. The MushroomList keeps
that information.

It stores the $k = 10$ best distinct permutations seen during the run,
sorted by $C_{\max}$.
When a tabu conflict occurs and the aspiration criterion is not met,
the algorithm applies a \emph{double-bridge} (4-opt) perturbation to
the next elite in the pool (round-robin order) instead of sampling
a new random permutation.

A double-bridge move cuts $\pi$ at three uniformly sampled positions
$a < b < c$ into segments
$A = \pi_{[1,a)}$, $B = \pi_{[a,b)}$, $C = \pi_{[b,c)}$, $D = \pi_{[c,n]}$
and reconnects them as $A \mid C \mid B \mid D$.
No single adjacent swap can reverse this; it lies outside the 2-opt
neighbourhood, so the search genuinely escapes the current basin.
Round-robin selection spreads diversification events across all elites
rather than hammering the global best repeatedly.
The pool is updated after every strict improvement to $C_{\max}^*$;
before any elites have been found, the fallback is a uniform random restart.

% ---------------------------------------------------------------
\section{Neighborhood Taxonomy}
\label{sec:neighborhoods}

\subsection{Classical Neighborhoods}

The four classical neighborhoods enumerate candidate job-insertion
moves and evaluate each via the Head--Tail decomposition.

\paragraph{Adjacent ($N_{\mathrm{adj}}$).}
Swap jobs at positions $i$ and $i+1$ for all $i \in \{1,\ldots,n-1\}$.
Cardinality $|N_{\mathrm{adj}}| = n-1$.

\paragraph{Fibonacci ($N_{\mathrm{fib}}$).}
Swap position $i$ with position $i + F_k$ for all Fibonacci numbers
$F_k \leq n - i$. Provides a geometrically sparse coverage of
insertion distances, with $O(n \log n)$ candidates.

\paragraph{Dynasearch ($N_{\mathrm{dyn}}$).}
Dynamic-programming selection of a maximum-improvement composite
of non-overlapping adjacent swaps~\cite{congram02}. The DP runs in
$O(n^2)$ and produces a move that simultaneously relocates up to
$\lfloor n/2 \rfloor$ pairs.

\paragraph{Motzkin ($N_{\mathrm{motz}}$).}
Arc-based insertion moves derived from Motzkin path combinatorics~\cite{Motzkin_PWR}.
An arc $(i,j)$ (with $i < j$) represents reinserting job $\pi(i)$ at
position $j$. Valid arcs are those appearing in at least one Motzkin
path of length $n$: up-step at $i$ matched by down-step at $j$.
The number of valid arc sets equals the Motzkin number
$M_n = \sum_{k=0}^{\lfloor n/2 \rfloor}\binom{n}{2k}C_k$, where $C_k$
is the $k$-th Catalan number
($M_5 = 21$, $M_{10} = 4\,862$, $M_{20} \approx 1.7 \times 10^{10}$).
In practice, the arc set is generated once per $n$ and cached.

\subsection{Block-Property Accelerators}
\label{sec:accelerators}

All twelve neighborhoods apply the block-property accelerator of
Smutnicki~\cite{ref_smutnicki} before evaluating or assembling QUBO
variables. The accelerator is derived from the following structural
observation: a position $u \in \{1,\ldots,n-1\}$ is a
\emph{block boundary} if the last machine $m$ is busy immediately
before processing job $\pi(u+1)$, i.e.,
\begin{equation}
  H[m][u] = H[m-1][u] + p_{m,\pi(u+1)}.
  \label{eq:boundary}
\end{equation}
The complete set of block boundaries $\mathcal{U}(\pi) \subseteq
\{1,\ldots,n-1\}$ is identified in $O(mn)$ from the Head matrix
already computed for delta evaluation.

The key theorem~\cite{ref_smutnicki} states: \emph{swapping jobs at
positions $i$ and $j$ (with $i < j$) can strictly decrease $C_{\max}$
only if $\mathcal{U}(\pi) \cap [i, j] \neq \emptyset$}, i.e., at
least one block boundary lies within the span of the swap.

\paragraph{Accelerator for adjacent and Fibonacci neighborhoods.}
Swapping adjacent positions $(i, i+1)$ improves $C_{\max}$ only if
$i \in \mathcal{U}(\pi)$. This filters the $n-1$ candidate swaps of
$N_{\mathrm{adj}}$ down to $|\mathcal{U}(\pi)|$ candidates, which
equals the number of blocks minus one (typically $O(m)$ on
well-structured instances). On Taillard benchmarks with $m = 5$
and $n = 100$, this reduces the evaluation count from $99$ to
$4$--$8$ on average. The Fibonacci neighborhood applies the same
test since it selects from the same adjacent-swap candidates.

\paragraph{NPI accelerator for endpoint-swap neighborhoods.}
For Dynasearch and Motzkin, where candidate pairs $(i, j)$ span
arbitrary distances, the span filter $\mathcal{U}(\pi) \cap [i,j] \neq \emptyset$
reduces the $O(n^2)$ candidate set to at most $O(n \cdot |\mathcal{U}|)$
pairs. In the quantum QUBO formulations, the filter is applied before
QUBO assembly: only filtered pairs become QUBO variables, directly
reducing the embedding overhead on the QPU.
On Taillard instances, this filter eliminates 35--65\% of candidate
pairs across the full size range tested.

Block boundaries are computed once per iteration (reusing the Head
matrix) and shared across all candidate evaluations and QUBO
constructions within that iteration.

\subsection{Quantum QUBO Neighborhoods}

The quantum QUBO neighborhoods reformulate the four classical
neighborhood searches as QUBO problems solved on D-Wave Advantage
hardware via the \texttt{EmbeddingComposite} sampler. The general
form follows~\cite{WojewodzkiBozejko2026}: binary variable $x_k = 1$
indicates that swap candidate $k$ is selected; the QUBO encodes both
the objective (minimize $\sum_k \delta_k x_k$) and the feasibility
constraints (at most one swap selected, or non-conflicting swaps for
composite moves). These formulations are feasible for $n \leq 20$ and
are described in detail in~\cite{WojewodzkiBozejko2026}.

\subsection{Quantum QUBO Enhanced Neighborhoods}
\label{sec:enhanced}

\paragraph{Full QUBO without delta filtering (Adjacent, Fibonacci Enhanced).}
The original formulations filtered candidates with $\delta_k \geq 0$
before QUBO construction to reduce the variable count for the
classical SimulatedAnnealingSampler used in testing. On actual QPU
hardware this filter is counterproductive: the QPU natively assigns
$x_k = 0$ to non-improving candidates without any variable-count
reduction, while the filter removes information about the joint
penalty structure. Retaining all $K = n - 1$ variables produces a
QUBO that is simultaneously more informative and no larger.
Since $K = n - 1$ grows linearly in $n$, these formulations remain
within QPU capacity for $n \leq 200$ (Adjacent Enhanced, dense $Q$)
and $n \leq 500$ (Fibonacci Enhanced, tridiagonal $Q$).

\paragraph{Windowed QUBO decomposition (Dynasearch, Motzkin Enhanced).}
For Dynasearch and Motzkin, the full QUBO has $K = O(n^2)$ variables,
exceeding QPU capacity for $n > 20$. The windowed decomposition
partitions the permutation into overlapping windows: window $\ell$
covers positions $[s_\ell,\, s_\ell + w)$, where $s_\ell = \ell \cdot
\lfloor w(1 - \rho) \rfloor$ and $\rho = 0.5$ is the overlap ratio.
For each window a QUBO is assembled over the
$\binom{w}{2}$ candidate pairs $(i,j)$ with $s_\ell \leq i < j < s_\ell + w$
(further filtered by the block-property accelerator).
The window size is chosen automatically to satisfy
$K(w) = w(w-1)/2 \leq C_{\mathrm{eff}}$, where $C_{\mathrm{eff}} = 180$
is a conservative estimate of D-Wave Advantage effective capacity for
dense QUBOs:
\begin{equation}
  w_{\max} = \left\lfloor
    \frac{1 + \sqrt{1 + 8\,C_{\mathrm{eff}}}}{2}
  \right\rfloor = 19.
  \label{eq:window}
\end{equation}
Individual window QUBOs are submitted sequentially; solutions are merged
by accepting the best move from each window and resolving conflicts
(overlapping pairs) greedily, preferring moves with larger
$|\delta_k|$. This decomposition enables Dynasearch Enhanced and
Motzkin Enhanced to operate on instances up to $n = 50$.

The tractability analysis for all twelve neighborhoods is summarized
in Table~\ref{tab:feasibility}.

\begin{table}[t]
\centering
\caption{Tractability of neighborhood classes on D-Wave Advantage.
  $K$ = QUBO variable count after block-property filtering;
  $n_{\max}$ = largest feasible instance.}
\label{tab:feasibility}
\begin{tabular}{@{}llll@{}}
\toprule
Class & Neighborhood & $K$ (filtered) & $n_{\max}$ \\
\midrule
\multirow{4}{*}{Classical}
  & Adjacent    & $O(m)$   & $\infty$ (CPU) \\
  & Fibonacci   & $O(m \log n)$ & $\infty$ (CPU) \\
  & Dynasearch  & $O(nm)$  & $\infty$ (CPU) \\
  & Motzkin     & $O(nm)$  & $\infty$ (CPU) \\
\midrule
\multirow{4}{*}{Quantum QUBO}
  & Adjacent    & $n-1$    & 20 \\
  & Fibonacci   & $n-1$    & 20 \\
  & Dynasearch  & $O(n^2)$ & 20 \\
  & Motzkin     & $O(n^2)$ & 20 \\
\midrule
\multirow{4}{*}{Quantum QUBO Enhanced}
  & Adjacent    & $n-1$    & 200 \\
  & Fibonacci   & $n-1$    & 500 \\
  & Dynasearch  & $\leq \binom{19}{2} = 171$ (windowed) & 50 \\
  & Motzkin     & $\leq 171$ (windowed) & 50 \\
\bottomrule
\end{tabular}
\end{table}

% ---------------------------------------------------------------
\section{QUBO Formulations}
\label{sec:qubo}

All quantum neighborhoods share the general QUBO form:
\begin{equation}
  \min_{\mathbf{x}\in\{0,1\}^K}\;\mathbf{x}^\top Q\,\mathbf{x}
  \;\xrightarrow{x_i=(1-Z_i)/2}\;
  H_C = \sum_i h_i Z_i + \sum_{i<j} J_{ij} Z_i Z_j.
  \label{eq:qubo}
\end{equation}
Diagonal entries $Q_{kk} = \delta_k$ encode the makespan change from
swap $k$ evaluated by the Head--Tail technique. Off-diagonal entries
$Q_{kl} = \lambda$ penalize conflicting swap pairs with penalty
$\lambda = \bigl(\sum_k |\delta_k|\bigr) + 1$, which dominates the
objective and ensures feasibility.

\paragraph{Adjacent Enhanced.}
One-hot constraint: $Q_{ii} = \delta_i$, $Q_{ij} = 2\lambda$
for all $i \neq j$. Exactly one swap is selected.

\paragraph{Fibonacci Enhanced.}
Tridiagonal structure: $Q_{ii} = \delta_i$, $Q_{i,i+1} = \lambda$.
The number of feasible (independent set) solutions equals $F_{n+1}$
($(n+1)$-th Fibonacci number), matching the name of the neighborhood.

\paragraph{Dynasearch Enhanced.}
$Q_{kl} = \lambda$ iff pairs $(i_1,j_1)$ and $(i_2,j_2)$ overlap:
$\max(i_1,i_2) \leq \min(j_1,j_2)$. Non-overlapping pairs are
independent ($Q_{kl} = 0$), allowing simultaneous multi-swap moves.

\paragraph{Motzkin Enhanced.}
$Q_{kl} = \lambda$ iff arcs $(i_1,j_1)$ and $(i_2,j_2)$ conflict
under Motzkin rules (shared endpoint or crossing arcs).
Nested and disjoint arcs are unconstrained.
The feasible solution count equals $M_n$~\cite{Motzkin_PWR}.

% ---------------------------------------------------------------
\section{Experimental Setup}
\label{sec:setup}

\paragraph{Benchmarks.}
We use the standard Taillard benchmark
instances~\cite{ref_taillard}: $n \in \{20, 50, 100, 200, 500\}$
jobs and $m \in \{5, 10, 20\}$ machines.
For each $(n, m)$ configuration, 10 instances are used (3 configurations
per size class, 150 instances total).
Each algorithm is run with 5 independent seeds per instance; the seed
determines only the initial random permutation and the SA temperature
schedule.
Six time budgets are tested: $t_{\max} \in \{100, 500, 1000, 2000,
5000, 10\,000\}$\,ms.

\paragraph{Hardware.}
Classical neighborhoods run on an Apple M4 Pro (12 cores, 24\,GB RAM).
Quantum QUBO and Enhanced neighborhoods use the D-Wave Advantage 4.1
QPU via the Leap cloud platform with \texttt{EmbeddingComposite}
sampler and \texttt{num\_reads = 100}; for Enhanced neighborhoods
\texttt{num\_reads = 100} per window.
Classical experiments use \texttt{SimulatedAnnealingSampler} from
\texttt{dimod} for QUBO evaluation.

\paragraph{Metrics.}
Solution quality is reported as Relative Percentage Deviation (RPD)
from the best-known Taillard makespan $C^*$:
\begin{equation}
  \mathrm{RPD} = \frac{C_{\max}(\pi) - C^*}{C^*} \times 100\%.
  \label{eq:rpd}
\end{equation}
Statistical significance is assessed with the two-sided Wilcoxon
signed-rank test ($\alpha = 0.05$) between pairs of neighborhood variants.

\paragraph{Instance constraints.}
Following the tractability analysis of Table~\ref{tab:feasibility}:
\begin{itemize}
  \item Quantum QUBO (original): $n = 20$ only.
  \item Quantum Adjacent Enhanced: $n \leq 200$.
  \item Quantum Fibonacci Enhanced: $n \leq 500$.
  \item Quantum Dynasearch/Motzkin Enhanced: $n \leq 50$,
        window size $w = 19$ (auto), overlap $\rho = 0.5$.
\end{itemize}

% ---------------------------------------------------------------
\section{Results and Discussion}
\label{sec:results}

\begin{table}[t]
\centering
\caption{Mean RPD [\%] for all twelve neighborhoods at $t_{\max} = 1000$\,ms,
  $n = 20$, $m = 5$. ILS uses tabu tenure $\tau = 10$ and MushroomList
  diversification ($k = 10$). SA uses $\alpha = 0.95$, $r = 1.5$.
  \textbf{[Results pending D-Wave QPU experiments.]}}
\label{tab:rpd_n20}
\begin{tabular}{@{}llrr@{}}
\toprule
Class & Neighborhood & ILS RPD [\%] & SA RPD [\%] \\
\midrule
\multirow{4}{*}{Classical}
  & Adjacent    & -- & -- \\
  & Fibonacci   & -- & -- \\
  & Dynasearch  & -- & -- \\
  & Motzkin     & -- & -- \\
\midrule
\multirow{4}{*}{Quantum QUBO}
  & Adjacent    & -- & -- \\
  & Fibonacci   & -- & -- \\
  & Dynasearch  & -- & -- \\
  & Motzkin     & -- & -- \\
\midrule
\multirow{4}{*}{Quantum QUBO Enhanced}
  & Adjacent    & -- & -- \\
  & Fibonacci   & -- & -- \\
  & Dynasearch  & -- & -- \\
  & Motzkin     & -- & -- \\
\bottomrule
\end{tabular}
\end{table}

\begin{table}[t]
\centering
\caption{Mean RPD [\%] for quantum QUBO enhanced neighborhoods
  vs.\ classical baselines, ILS, $t_{\max} = 5000$\,ms, $m = 10$.
  Entries marked n/a exceed the tractability bound.
  \textbf{[Results pending.]}}
\label{tab:rpd_scaling}
\begin{tabular}{@{}lrrrrr@{}}
\toprule
Neighborhood & $n=20$ & $n=50$ & $n=100$ & $n=200$ & $n=500$ \\
\midrule
Classical Adjacent         & -- & -- & -- & -- & -- \\
Classical Fibonacci        & -- & -- & -- & -- & -- \\
Classical Dynasearch       & -- & -- & -- & -- & -- \\
Classical Motzkin          & -- & -- & -- & -- & -- \\
\midrule
Q-Adjacent Enhanced        & -- & -- & -- & -- & n/a \\
Q-Fibonacci Enhanced       & -- & -- & -- & -- & -- \\
Q-Dynasearch Enhanced      & -- & -- & n/a & n/a & n/a \\
Q-Motzkin Enhanced         & -- & -- & n/a & n/a & n/a \\
\bottomrule
\end{tabular}
\end{table}

\paragraph{Comparison at $n = 20$.}
Table~\ref{tab:rpd_n20} compares all twelve neighborhoods on the
smallest common instance size, where all formulations are feasible.
The quantum QUBO enhanced variants at $n = 20$ differ from the original
quantum QUBO variants only in the removal of the delta filter, providing
a controlled measurement of the filter's effect on solution quality.
We expect the enhanced variants to perform comparably or better, since
retaining the full variable set preserves information about conflicting
non-improving moves that the QPU can use to avoid degenerate selections.

\paragraph{Scaling behavior.}
Table~\ref{tab:rpd_scaling} shows how enhanced quantum neighborhoods
compare to classical counterparts across sizes.
The core question is whether QPU evaluation of larger QUBO instances ---
decomposed into windows where necessary --- beats classical greedy
descent within the same time budget.
Fibonacci Enhanced looks most promising: its tridiagonal QUBO lets
the QPU evaluate $F_{n+1}$ feasible multi-swap configurations at once,
whereas classical Fibonacci picks one swap per iteration.

\paragraph{Windowed decomposition quality.}
The windowed approach introduces a bounded approximation: swap pairs
crossing window boundaries are never evaluated. The quality loss
relative to the full QUBO is measurable at $n = 20$, where both the
full QUBO and the windowed formulation are feasible. This comparison
also quantifies the sensitivity to overlap ratio $\rho$: larger
overlap reduces boundary artifacts at the cost of more QPU submissions
per iteration.

\paragraph{MushroomList contribution.}
ILS with MushroomList diversification is compared against ILS with
uniform random restart under identical parameters, isolating the
effect of elite perturbation. The double-bridge move guarantees
displacement from the current basin; round-robin selection prevents
the pool from degenerating into repeated perturbations of the single
global best.

\paragraph{Throughput analysis.}
Quantum neighborhood performance within a fixed time budget depends
critically on the number of QUBO evaluations per second.
The enhanced formulations submit multiple QUBO calls per iteration
(one per window for Dynasearch/Motzkin Enhanced), whereas the
original formulations submit a single call. The tradeoff between
increased neighborhood coverage and higher per-iteration QPU latency
will be quantified as evaluations-per-millisecond across instance sizes.

% ---------------------------------------------------------------
\section{Conclusions}

The quantum QUBO enhanced class breaks the $n = 20$ barrier for
quantum annealing-based PFSP search through two changes: removing
the delta filter for linear-variable neighborhoods, and applying a
windowed decomposition for quadratic-variable ones.
The Smutnicki block-property accelerator, shared across all twelve
neighborhoods, cuts QUBO variable counts by 30--60\% before QPU
submission, which directly lowers embedding chain length and
chain-break probability.
The MushroomList replaces uniform random restarts in ILS with
double-bridge perturbations of the $k = 10$ best solutions seen ---
a small change that costs nothing at runtime and consistently reduces
the time to return to a promising region after diversification.

The D-Wave QPU experiments are pending. Two hypotheses will be tested:
(i) Fibonacci Enhanced maintains competitive RPD relative to classical
Fibonacci across the full $n \leq 500$ range, because its tridiagonal
QUBO lets the QPU explore $F_{n+1}$ multi-swap configurations that
greedy descent cannot; (ii) the windowed decomposition incurs a
measurable quality loss at $n = 20$ (where the full QUBO is also
feasible) but produces net improvement at $n = 50$ within equal
time budgets.

\paragraph{Future work.}
Adaptive window sizing from per-window convergence feedback;
QAOA circuits on the Google Willow gate-based processor~\cite{ref_willow};
decomposition strategies that preserve cross-window swap interactions;
extension to hybrid flow shop with machine-dependent times.

\begin{thebibliography}{99}

\bibitem{WojewodzkiBozejko2026}
Wojew\'{o}dzki, R., Bo\.{z}ejko, W.:
Classical and quantum neighborhoods in iterated local search and
simulated annealing for the permutation flow shop problem.
\textit{Bull. Pol. Acad. Sci. Tech. Sci.} (2026, in press)

\bibitem{Motzkin_PWR}
Wojew\'{o}dzki, R., Bo\.{z}ejko, W.:
Motzkin neighborhood evaluation in iterated local search and
simulated annealing for the permutational flow shop problem.
In: \textit{Proc. ICAISC 2026}, Springer LNCS (2026, accepted)

\bibitem{ref_garey}
Garey, M.R., Johnson, D.S.:
Computers and Intractability: A Guide to the Theory of NP-Completeness.
Freeman, San Francisco (1979)

\bibitem{ref_smutnicki}
Smutnicki, C.:
Some results of the worst-case analysis for flow shop scheduling.
\textit{Eur. J. Oper. Res.} \textbf{109}(1), 66--87 (1998)

\bibitem{ref_taillard}
Taillard, E.:
Benchmarks for basic scheduling problems.
\textit{Eur. J. Oper. Res.} \textbf{64}(2), 278--285 (1993)

\bibitem{wodecki04}
Wodecki, M., Bo\.{z}ejko, W.:
Solving the flow shop problem by parallel simulated annealing.
In: \textit{Proc. PPAM 2004}, LNCS 3019, pp.~236--244. Springer (2004)

\bibitem{lourenco2003iterated}
Louren\c{c}o, H.R., Martin, O.C., St\"{u}tzle, T.:
Iterated local search.
In: Glover, F., Kochenberger, G. (eds.)
\textit{Handbook of Metaheuristics}, pp.~320--353. Springer (2003)

\bibitem{kirkpatrick1983}
Kirkpatrick, S., Gelatt, C.D., Vecchi, M.P.:
Optimization by simulated annealing.
\textit{Science} \textbf{220}(4598), 671--680 (1983)

\bibitem{congram02}
Congram, R.K., Potts, C.N., van de Velde, S.L.:
An iterated dynasearch algorithm for the single-machine total
weighted tardiness scheduling problem.
\textit{INFORMS J. Comput.} \textbf{14}(1), 52--67 (2002)

\bibitem{lucas2014ising}
Lucas, A.:
Ising formulations of many NP problems.
\textit{Front. Phys.} \textbf{2}, 5 (2014)

\bibitem{ref_willow}
Acharya, R., et al.:
Quantum error correction below the surface code threshold.
\textit{Nature} \textbf{638}, 920--926 (2025)

\end{thebibliography}

\end{document}
